% !TeX root = ../../Book2.tex
% 纲要标题：## 第5章　局部系统与扭系数同调
% 纲要位置：几何学计划纲要.md，第 1191 行。
% 正文状态：初稿；尚未正式审核。

% 纲要内容开始
% 
% > 本章目标与证明范围：从覆盖的单值进入非平凡系数，完整构造链模型和函子性。定向局部系统、纤维化中的系数作用及障碍系数在此有共同的实际基础。 分类主版本固定连通、局部道路连通和半局部单连通空间；更一般群胚表述另述。
% 
% 纲要内容结束

\chapter{局部系统与扭系数同调}
\label{g2:ch:05}
\BookChapterNumber{g2:ch:05}

\begin{chapterabstract}
	本章把固定系数推广为随位置变化、可沿道路运输的系数。先建立路径群胚、局部常值层与基本群表示之间的联系，再用万有覆盖上的群环链复形构造扭系数同调及上同调，并给出首顶点系数的几何公式。定向局部系统说明符号单值如何来自流形自身；圆周和不可定向曲面的计算展示核、余核及顶次整数类怎样随单值改变。最后把纤维同调组织为基底上的局部系统，证明映射环面的正合列，为后续纤维化、谱序列与障碍理论准备系数数据。
\end{chapterabstract}

\begin{learninggoals}
	\begin{itemize}
		\item 用道路运输定义局部系统，在明确的局部条件下证明它与局部常值层、基本群模的等价，说明换基点时需要保留的比较映射。
		\item 区分左右群环作用，构造局部系数链与上链，证明函子性、正合列、同伦不变性、切除及胞腔计算的依据。
		\item 从流形局部同调构造定向系统，计算圆周与不可定向曲面的扭系数同调及上同调，解释不变元和余不变元的差别。
		\item 构造纤维同调与上同调的运输，计算映射环面，并区分局部系统的不变类、几何截面和更高延拓问题。
	\end{itemize}
\end{learninggoals}

在前两章的链中，一个单形的系数是整数，取边界时把同一个整数放在各个面前即可。这个习惯默认了不同位置的系数已经处在同一个群中。可是，在不可定向曲面上移动一小片有定向的圆盘，绕某条闭路回来时，其局部定向可能反转。若系数用来记录这片圆盘的定向，那么“绕一圈后仍取原系数”就不再符合几何。

\ProseParagraphBreak

另一种情形来自一族空间。各个纤维的同调群可能彼此同构，但沿基底环路比较它们以后，这个群会经历一个非平凡自同构。只记录“每个纤维的群都是 \(A\)”会漏掉这个自同构；局部系统把群和比较它们的道路运输一起保留下来。上一章改变的是所有位置共用的系数，本章还允许绕行改变系数的识别。

\ProseParagraphBreak

\chapref{g2:ch:01}的覆盖与单值给出运输的几何模型，\chapref{g2:ch:02}的奇异链和正合列用于构造带运输的边界，\chapref{g2:ch:03}的胞腔方法则将计算缩到少量生成元。上链及其乘积沿用\chapref{g2:ch:04}的构造。层在这里记录可以局部选择并粘合的系数，其定义和与运输的联系在本章建立。路径群胚上的运输可在任意空间上定义；它与基本群表示、局部常值层之间的分类对应，在连通、局部道路连通、半局部单连通的条件下证明。

\ProseParagraphBreak

本章的一个持续约定是：\(\alpha*\beta\) 表示先走 \(\alpha\)、再走 \(\beta\)，故运输满足 \(T_{\alpha*\beta}=T_\beta T_\alpha\)。它决定后面的右群作用、首顶点边界和上链中的逆映射。局部系数的阅读对照主要见\cite[第3.H节]{Hatcher2002}；不同书的左右作用与路径复合约定可能不同，比较公式时须同时比较这些约定。

% !TeX root = ../../Book2.tex
% 纲要标题：### 5.1 局部系统的等价描述
% 纲要位置：几何学计划纲要.md，第 1195 行。
% 正文状态：初稿；尚未正式审核。

% 纲要内容开始
% 
% * 路径群胚上的运输函子
% * 局部常值层与纤维
% * 连通局部良好空间上的基本群表示
% * 换基点和同构的相容
% 
% 纲要内容结束

\section{局部系统的等价描述}
\label{g2:sec:0501}

在通常的整数链中，单形上的系数可以直接相加，因为它们都属于同一个指定的群 \(\mathbb Z\)。如果每一点只给出一个同构于 \(\mathbb Z\) 的群，情形就不同了：把一个系数送到另一点需要选择道路，绕一圈回来还可能把它变成相反数。定向的改变正会产生这种现象。本节先把“沿道路比较系数”本身定义清楚，再说明在适当的局部条件下，它也可以表现为一个局部常值的层，或一个基本群作用。

\ProseParagraphBreak

以下 \(R\) 为含幺交换环，系数纤维为 \(R\) 模，\(R=\mathbb Z\) 时就是阿贝尔群。这里不要求模自由或有限生成。沿用\cref{g2:def:fundamental-group}的顺序，\(\alpha*\beta\) 表示先走 \(\alpha\)、后走 \(\beta\)，\(\bar\alpha\) 为反向道路；因此运输的复合写成 \(T_{\alpha*\beta}=T_\beta\circ T_\alpha\)。

\subsection{路径群胚上的运输函子}
\label{g2:subsec:0501:01}
同一端点的环路组成基本群；若同时保留所有点以及不同点之间的道路，就得到一个能直接描述运输的对象。

\begin{definition}[局部系统]\label{g2:def:local-system}
	设 \(X\) 为拓扑空间。它的\term[lujingqunpei]{路径群胚}[fundamental groupoid] \(\Pi_1(X)\) 以 \(X\) 的点为对象，以定端点同伦道路类 \([\alpha]:x\to y\) 为态射，复合为
	\[
		[\beta]\circ[\alpha]=[\alpha*\beta].
	\]
	常值道路给出恒等态射，反向道路给出逆态射；这些运算良定义由\cref{g2:def:fundamental-group}之后对重参数化和往返道路的论证保证。\symidx{\(\Pi_1(X)\)：路径群胚}

	一个 \(R\) 模\term[jubuxitong]{局部系统}[local system] \(\mathcal L\) 是从 \(\Pi_1(X)\) 到 \(R\) 模范畴的函子。具体说，它给每点 \(x\) 一个模 \(\mathcal L_x\)，给每条道路类一个线性同构 \(T_\alpha:\mathcal L_x\to\mathcal L_y\)，并满足
	\begin{equation}\label{g2:eq:local-transport-composition}
		T_{e_x}=\operatorname{id},\qquad T_{\alpha*\beta}=T_\beta T_\alpha.
	\end{equation}
	\(\mathcal L_x\) 称为在 \(x\) 处的纤维，\(T_\alpha\) 称为沿 \(\alpha\) 的运输。逆道路恒等式给出 \(T_{\bar\alpha}=T_\alpha^{-1}\)。局部系统之间的态射 \(u:\mathcal L\to\mathcal K\) 是一族线性映射 \(u_x:\mathcal L_x\to\mathcal K_x\)，满足每条 \(\alpha:x\to y\) 上的自然性
	\[
		\begin{tikzcd}[column sep=large,row sep=large]
			\mathcal L_x \arrow[r,"T^{\mathcal L}_\alpha"] \arrow[d,"u_x"'] & \mathcal L_y \arrow[d,"u_y"] \\
			\mathcal K_x \arrow[r,"T^{\mathcal K}_\alpha"'] & \mathcal K_y.
		\end{tikzcd}
	\]
	\symidx{\(\mathcal L_x,T_\alpha\)：局部系统的纤维与道路运输}
\end{definition}

这个定义在任意拓扑空间上有意义；它只使用道路及其同伦，没有预先给纤维的并集指定拓扑。两个局部系统同构，指存在各纤维上均为同构的态射，其逆映射族自动满足同一自然性方形。固定一个 \(R\) 模 \(M\)，在每点都取 \(M\)、每条道路都取恒等映射，得到\term[changzhixitong]{常值系统}[constant system]，记为 \(\underline M\)。局部系统平凡，是指它同构于某个这样的系统；这并不表示已经选定那个同构。

\begin{proposition}[拉回与同伦运输]\label{g2:prop:local-system-pullback}
	设 \(f:Y\to X\) 为连续映射，\(\mathcal L\) 为 \(X\) 上的 \(R\) 模局部系统。规定
	\[
		(f^*\mathcal L)_y=\mathcal L_{f(y)},\qquad T^{f^*\mathcal L}_\alpha=T^{\mathcal L}_{f\circ\alpha},
	\]
	便得到 \(Y\) 上的局部系统，称为 \(\mathcal L\) 的\term[lahuijubuxitong]{拉回局部系统}[pullback local system]。它也逐纤维拉回态射，并满足 \((f\circ g)^*=g^*f^*\) 及恒等映射的相容性。若 \(H:Y\times I\to X\) 是从 \(f_0\) 到 \(f_1\) 的同伦，令 \(\lambda_y(t)=H(y,t)\)，则
	\[
		T_{\lambda_y}:(f_0^*\mathcal L)_y\longrightarrow(f_1^*\mathcal L)_y
	\]
	组成局部系统同构。此同构由所给同伦确定；不同的同伦可能给出不同的同构。
\end{proposition}
\begin{proof}
	连续映射把定端点同伦送到定端点同伦，并保持道路拼接及常值道路，所以拉回满足\eqref{g2:eq:local-transport-composition}。态射、恒等和复合的公式都是对相应点或道路逐项求值。

	对道路 \(\alpha:y\to z\)，将 \(H\) 限制在 \(\alpha\) 上得到方形 \((s,t)\mapsto H(\alpha(s),t)\)。方形内两条从左下到右上的边界路线定端点同伦，故
	\[
		(f_0\alpha)*\lambda_z\simeq\lambda_y*(f_1\alpha).
	\]
	应用运输公式即得 \(T_{\lambda_z}T_{f_0\alpha}=T_{f_1\alpha}T_{\lambda_y}\)，这正是态射的自然性。每个 \(T_{\lambda_y}\) 可逆，故得到同构。取 \(Y\) 为一点、\(f_0=f_1\) 时，同伦就是基点环路，所得同构就是该环路的单值；因此它确实可能依赖同伦。
\end{proof}

拉回的定义解释了系数为何需要跟着空间映射一起变化：若 \(y\) 映到 \(f(y)\)，在 \(y\) 使用的系数正应取自 \(\mathcal L_{f(y)}\)。稍后定义链映射时，这会成为系数所在位置的相容条件。

\subsection{局部常值层与纤维}
\label{g2:subsec:0501:02}
道路运输给出不同纤维的比较，层则记录一个开集上能同时选取哪些彼此相容的系数。在运输平凡的小邻域内，选定一个值后可把它搬到各点，得到一份局部常值截面；层的限制与粘合正是为了组织这样的局部选择。

\begin{definition}[局部常值层]\label{g2:def:locally-constant-sheaf}
	拓扑空间 \(X\) 上的 \(R\) 模\term[ceng]{层}[sheaf] \(\mathscr F\)，给每个开集 \(U\) 一个 \(R\) 模 \(\mathscr F(U)\)，并给每个包含 \(V\subset U\) 一个线性限制映射 \(s\mapsto s|_V\)。限制保持恒等与复合，且满足粘合条件：若 \(U=\bigcup_iU_i\)，而 \(s_i\in\mathscr F(U_i)\) 在每个交集上相等，则存在唯一 \(s\in\mathscr F(U)\) 使 \(s|_{U_i}=s_i\)。\(\mathscr F(U)\) 的元素称为 \(U\) 上的\term[jiemian]{截面}[section]。

	在 \(x\) 处，把定义在 \(x\) 的某个开邻域上的截面按“在更小邻域上一致”识别，所得模称为\term[jing]{茎}[stalk] \(\mathscr F_x\)；截面 \(s\) 的类记为 \(s_x\)，称为它在 \(x\) 处的\term[ya]{芽}[germ]。加法和数乘在两个代表的公共邻域内进行，进一步缩小邻域不改变结果，故定义良好。

	给模 \(A\) 离散拓扑，\(U\mapsto\{\text{局部常值函数 }U\to A\}\) 构成常值层 \(\underline A\)。若每个 \(x\in X\) 都有开邻域 \(U\)，使限制在 \(U\) 上的层 \(\mathscr F|_U\) 同构于 \(\underline A|_U\)，其中 \(A\) 可随邻域变化，则称 \(\mathscr F\) 为\term[jubuchangzhiceng]{局部常值层}[locally constant sheaf]。层的态射是一族与限制相容的线性映射。
\end{definition}

这里的常值层允许一个不连通开集的不同部分取不同值。若仅给每个非空开集同一个模 \(A\)，并把所有限制设为恒等，就通常不满足粘合条件：两个不交开集上选的不同常值无法粘成一个共同常值。另一方面，茎只看一点附近，所以 \(\underline A_x\cong A\)，同构把一个局部常值函数的芽送到它在 \(x\) 的值。

\begin{proposition}[层的纤维空间]\label{g2:prop:local-sheaf-fiber-space}
	设 \(X\) 为拓扑空间，\(R\) 为含幺交换环。局部常值 \(R\) 模层等价于下列数据：一个覆盖 \(p:E\to X\)，各纤维具有 \(R\) 模结构，并在每点附近有保持纤维线性结构的同胚 \(p^{-1}(U)\cong U\times A\)，其中 \(A\) 取离散拓扑。态射分别取层态射和与投影相容、在各纤维上线性的连续映射。这样的 \(E\) 通过道路提升确定一个 \(R\) 模局部系统。
\end{proposition}
\begin{proof}
	先由层构造空间。令 \(E(\mathscr F)=\coprod_{x\in X}\mathscr F_x\)，投影把芽送到它所在的点。对 \(s\in\mathscr F(U)\)，定义 \(\widehat s(U)=\{s_x:x\in U\}\)。以这些集合为开集的基：若 \(s_x=t_x\)，芽的定义给出 \(x\) 的邻域 \(W\)，使 \(s|_W=t|_W\)，所以 \(\widehat s(W)=\widehat t(W)\) 包含在原来两个集合的交内。每个芽有截面代表，故这些集合也覆盖 \(E(\mathscr F)\)。投影在 \(\widehat s(U)\) 上是到 \(U\) 的双射，其逆 \(x\mapsto s_x\) 连续；后一结论再次来自“两个芽相同就在某个邻域上一致”。投影本身连续，因为开集的逆像是这类基本开集的并。因此投影在每个 \(\widehat s(U)\) 上是同胚。

	若 \(\mathscr F|_U\cong\underline A|_U\)，茎的对应给出 \(E(\mathscr F)|_U\cong U\times A\)。在这个对应中，基本开集就是局部常值函数的图像；它们与离散 \(A\) 的乘积拓扑一致。因此投影不仅是局部同胚，而且在同一 \(U\) 上把全部纤维同时平凡化，确为覆盖。一个层态射在茎上诱导线性映射；它把截面芽族送到像截面的芽族，所以得到连续的纤维线性映射。

	反过来，对题述 \(E\) 定义 \(\mathscr S_E(U)\) 为所有连续截面 \(s:U\to E\)。加法和数乘逐点进行，在乘积邻域中就是对局部常值函数进行运算，因而仍连续。相容截面逐点粘合后连续，且唯一，故 \(\mathscr S_E\) 是层；在每个乘积邻域上，它正是常值层。映射 \(\mathscr F(U)\to\mathscr S_{E(\mathscr F)}(U)\)、\(s\mapsto(x\mapsto s_x)\) 单射，因为芽处处相同的截面局部相同，再由粘合唯一性相同。它满射，因为一个连续截面在任一点附近落入某个 \(\widehat t(V)\)，于是局部来自截面 \(t\)，这些代表在交集上具有相同的芽，因而能够唯一粘合。另一方面，\(E(\mathscr S_E)\to E\) 把截面的芽送到该点的值；在乘积邻域中它是恒等的纤维对应，故为纤维线性同胚。这两个对应也保持态射。

	最后，\cref{g2:thm:covering-lifts}已证明：覆盖中从指定纤维点出发的道路提升唯一，提升终点只依赖定端点同伦类。把提升终点定义为 \(T_\alpha\)，拼接提升给出\eqref{g2:eq:local-transport-composition}。将道路有限细分到乘积邻域，在每一段内运输都是常值坐标的识别，所以线性，复合后仍线性。连续纤维线性映射把道路提升送到道路提升，由唯一性便与运输交换。
\end{proof}

纤维空间使用离散的系数拓扑。即使 \(R=\mathbb R\)、纤维为 \(\mathbb R^r\)，这里也不是把通常的向量空间拓扑放在纤维上；离散拓扑使“沿道路保持局部常值”成为一项确定的数据。后文谈平坦丛时，可以把这项运输结构与其对应的光滑向量丛联系起来。

\begin{proposition}[运输的局部实现]\label{g2:prop:local-system-realization}
	设 \(X\) 局部道路连通且半局部单连通。每个 \(R\) 模局部系统 \(\mathcal L\) 的纤维并集 \(\coprod_x\mathcal L_x\) 都有自然的覆盖拓扑及局部模平凡化，使其道路提升恢复 \(\mathcal L\) 的运输。它与\cref{g2:prop:local-sheaf-fiber-space}的构造一起，给出局部系统与局部常值 \(R\) 模层的等价。
\end{proposition}
\begin{proof}
	由\cref{g2:def:local-covering-conditions}及\cref{g2:thm:universal-cover}证明开头的取邻域论证，可以取由道路连通开集 \(U\) 构成的开基，且 \(U\) 中的每条环路在 \(X\) 中零伦。固定 \(u\in U\)。若 \(y\in U\)，从 \(u\) 到 \(y\) 的两条 \(U\) 内道路的差环路在 \(X\) 中零伦，所以它们在 \(X\) 中定端点同伦，给出相同的运输。对每个 \(m\in\mathcal L_u\) 因而有确定的族
	\[
		s_m(y)=T_\eta(m),\qquad \eta:u\to y\text{ 位于 }U.
	\]
	对任意 \(y\)，\(m\mapsto s_m(y)\) 是线性同构，故这些族互不相交且填满 \(U\) 上的纤维。

	把全部集合 \(s_m(U)\) 宣布为开集的基。若两个这样的族在 \(y\) 处相同，在两个底开集的交内取含 \(y\) 的上述小邻域 \(W\)。两族在 \(W\) 上均由共同的值沿 \(W\) 内道路运输得到，故在整个 \(W\) 上相同，证明了基的交集条件。投影在每个 \(s_m(U)\) 上与 \(U\) 同胚，并有保持模结构的局部同胚
	\[
		U\times\mathcal L_u\longrightarrow p^{-1}(U),\qquad(y,m)\longmapsto s_m(y).
	\]
	这给出所求覆盖。道路可有限细分到这些开集；每一段的提升正沿某个 \(s_m\) 走，其终点就是原来的运输，拼接后便恢复整条道路的运输。

	局部系统态射在这样的邻域中由一点的线性映射决定，因此在所构造的覆盖上连续。若一开始已有局部模平凡化覆盖，再取同时包含在平凡化邻域内的上述小开集，则刚刚构造的 \(s_m\) 恰好是原覆盖的叶，故原拓扑被恢复。结合\cref{g2:prop:local-sheaf-fiber-space}对层与覆盖双向恢复的证明，得到连同态射在内的等价。
\end{proof}

半局部单连通性使小开集内的环路在整个空间中零伦，其运输因而为恒等。配合局部道路连通性，一个纤维中的值便能沿小开集内的道路唯一地展开为局部常值截面。这给出了路径运输与局部常值层在上述局部条件下的对应。

\subsection{连通局部良好空间上的基本群表示}
\label{g2:subsec:0501:03}
现在假设 \(X\) 非空、连通、局部道路连通且半局部单连通。连通和局部道路连通使 \(X\) 道路连通：道路分支都是开集，多于一个分支便会给出分离。任选 \(x_0\in X\)，从这里到每一点都能走道路，所以一个纤维加上环路运输理应包含整个局部系统的信息。需要说明的是，如何从这份信息重建其他纤维，并把道路选择的歧义准确地除掉。

\begin{theorem}[局部系统的分类]\label{g2:thm:local-system-classification}
	设 \(X\) 非空、连通、局部道路连通且半局部单连通，\(R\) 为含幺交换环，\(x_0\in X\)，\(\pi=\pi_1(X,x_0)\)。取基点纤维及其单值
	\begin{equation}\label{g2:eq:local-right-action}
		M=\mathcal L_{x_0},\qquad m\cdot g=T_g(m)
	\end{equation}
	给出以下三类数据之间的等价：\(X\) 上的 \(R\) 模局部系统、\(X\) 上的局部常值 \(R\) 模层、以及带 \(R\) 线性右 \(\pi\) 作用的 \(R\) 模。第三类数据的态射为保持右作用的 \(R\) 线性映射。

	具体地，对给定的右 \(\pi\) 模 \(M\)，取道路类万有覆盖 \(p:\widetilde X\to X\)，其中 \(g[\gamma]=[g*\gamma]\) 为左覆盖变换作用。相应的系数纤维空间为
	\begin{equation}\label{g2:eq:associated-local-system}
		E_M=(M\times\widetilde X)/{\sim},\qquad(m\cdot g,\widetilde x)\sim(m,g\widetilde x),
	\end{equation}
	\(M\) 取离散拓扑，\(E_M\) 取商拓扑，投影为 \([m,\widetilde x]\mapsto p(\widetilde x)\)。
\end{theorem}
\begin{proof}
	先检查右作用。由\eqref{g2:eq:local-transport-composition}，\((m\cdot g)\cdot h=T_hT_g(m)=T_{gh}(m)=m\cdot(gh)\)，而单位环路作用为恒等。局部系统态射在 \(x_0\) 处与每条环路的运输交换，所以得到保持此右作用的线性映射。

	反向构造用\cref{g2:thm:universal-cover,g2:prop:monodromy-deck}已证明的万有覆盖及左作用。式\eqref{g2:eq:associated-local-system}的等价关系正是左作用
	\[
		g\cdot(m,\widetilde x)=(m\cdot g^{-1},g\widetilde x)
	\]
	的轨道关系：先作用 \(h\) 再作用 \(g\)，第一坐标成为 \(m\cdot h^{-1}g^{-1}=m\cdot(gh)^{-1}\)，第二坐标成为 \(gh\widetilde x\)。每个作用映射是同胚，故轨道商映射是开映射，其开集饱和为这些同胚像的并。

	选万有覆盖的一个道路连通均匀覆盖邻域 \(U\) 及其上一叶，记相应的逆映射为 \(s:U\to\widetilde X\)。\(\pi\) 在每个万有覆盖纤维上自由且传递：对于同终点道路类 \([\alpha],[\beta]\)，唯一的 \(g=[\beta*\bar\alpha]\) 满足 \(g[\alpha]=[\beta]\)。覆盖变换因此把所选叶送到全部其他叶。每个 \(E_M\) 中底点位于 \(U\) 的元素，都唯一写成 \([m,s(y)]\)，因为 \([m,gs(y)]=[m\cdot g,s(y)]\)。于是
	\[
		U\times M\longrightarrow E_M|_U,\qquad(y,m)\longmapsto[m,s(y)]
	\]
	为双射。它连续，并且把每个 \(V\times\{m\}\)、\(V\subset U\) 开，送到开集：该集合在商前的饱和集是 \(\{m\cdot g^{-1}\}\times gs(V)\) 对全部 \(g\) 的并。故这是同胚。换一叶只把纤维坐标改为 \(m\mapsto m\cdot g\)，这是线性同构，所以每个纤维的模结构与选择无关。

	用 \(m\mapsto[m,\widetilde x_0]\)、\(\widetilde x_0=[e_{x_0}]\)，识别基点纤维与 \(M\)。代表 \(g\) 的环路从 \(\widetilde x_0\) 提升后终点为 \(g\widetilde x_0\)，因此 \(E_M\) 中相应运输的终值为
	\[
		[m,g\widetilde x_0]=[m\cdot g,\widetilde x_0],
	\]
	恰好恢复所给右作用。

	若 \(M\) 来自已有局部系统 \(\mathcal L\)，定义
	\[
		\Phi:E_M\longrightarrow\coprod_x\mathcal L_x,\qquad
		\Phi([m,[\gamma]])=T_\gamma(m).
	\]
	它尊重商关系，因为 \(T_{g*\gamma}(m)=T_\gamma(m\cdot g)\)。固定某条从 \(x_0\) 到 \(x\) 的道路，\(T_\gamma\) 给出每个纤维上的线性双射。在\cref{g2:prop:local-system-realization}的局部平凡化中，\(\Phi\) 是由固定线性同构给出的乘积映射，故连续，且逆连续；沿任一道路接上尾段也说明它保持运输。因此恢复的是原局部系统。

	最后，保持右作用的线性映射 \(u:M\to N\) 诱导 \([m,\widetilde x]\mapsto[u(m),\widetilde x]\)，由商关系可知良定义，并在局部平凡化中连续且线性。任意局部系统态射都由基点分量唯一确定，因为沿 \(\gamma:x_0\to x\) 的自然性要求
	\[
		u_x=T^{\mathcal K}_\gamma\,u_{x_0}\,(T^{\mathcal L}_\gamma)^{-1}.
	\]
	上述构造给出每一个保持右作用的基点映射，所以不只对象的同构类对应，态射也一一对应。局部常值层的等价已由\cref{g2:prop:local-system-realization}证明。
\end{proof}

按已经固定的道路拼接顺序，\(g\mapsto T_g\) 满足 \(T_{gh}=T_hT_g\)，因而给出右作用。取逆便将复合顺序转换为通常的左表示：
\begin{equation}\label{g2:eq:local-left-representation}
	\rho:\pi\longrightarrow\operatorname{Aut}_R(M),\qquad\rho(g)=T_g^{-1}.
\end{equation}
此时 \(\rho(gh)=\rho(g)\rho(h)\)。当 \(M=R^r\) 并选定基时，这成为到 \(\operatorname{GL}_r(R)\) 的表示；换纤维基使表示同时共轭，而非产生一个新的局部系统同构类。群丛、群环模与局部系数的关联还可参见 \cite[第3.H节]{Hatcher2002}；比较其中的公式时须一并转换左右作用。

\begin{example}[圆周上的符号系统]\label{g2:exm:sign-local-system}
	取 \(X=S^1\)、\(R=\mathbb Z\)、\(M=\mathbb Z\)。由\cref{g2:prop:circle-fundamental-group}，\(\pi_1(S^1,1)=\mathbb Z\)，正向生成元记为 \(t\)。令 \(m\cdot t=-m\)，得到一个局部系统 \(\mathbb Z_-\)。在 \(\widetilde X=\mathbb R\)、\(p(s)=\exp(2\pi\mathrm{i}s)\) 的模型中，它的纤维空间为
	\[
		(\mathbb Z\times\mathbb R)/\bigl((m,s+1)\sim(-m,s)\bigr).
	\]
	沿一小段圆弧，整数系数保持不变；绕圆一周以后，同一起点纤维内的整数变成相反数。它不能同构于 \(\underline{\mathbb Z}\)：若基点同构为 \(u\)，自然性将要求 \(u(-m)=u(m)\)，取 \(m=1\) 即与 \(u\) 单射矛盾。它的每个纤维仍然都是 \(\mathbb Z\)，可见只列出纤维的同构类型不够。
\end{example}

\subsection{换基点和同构的相容}
\label{g2:subsec:0501:04}
分类中选择基点，是为了用一个模保存全部纤维的信息。改用另一个基点不会改变局部系统，但比较两个模时必须给出连接道路。\cref{g2:prop:basepoint-change}已用往返道路的消去证明基本群的换基点公式，下面将它与系数运输同时写出。

\begin{proposition}[换基点与单值]\label{g2:prop:basepoint-local-system}
	设 \(\mathcal L\) 为拓扑空间 \(X\) 上的 \(R\) 模局部系统，\(c:x_0\to x_1\) 为道路，\(M_i=\mathcal L_{x_i}\)。定义
	\[
		\psi_c:\pi_1(X,x_1)\longrightarrow\pi_1(X,x_0),\qquad
		\psi_c([h])=[c*h*\bar c].
	\]
	这是群同构，且 \(T_c:M_0\to M_1\) 满足
	\begin{equation}\label{g2:eq:basepoint-module-transport}
		T_c(m\cdot\psi_c(h))=T_c(m)\cdot h.
	\end{equation}
	若 \(c':x_0\to x_1\) 为另一条道路，令 \(g=[c'*\bar c]\in\pi_1(X,x_0)\)，则
	\[
		T_{c'}=T_cT_g,\qquad\psi_{c'}(h)=g\psi_c(h)g^{-1}.
	\]
	局部系统态射 \(u:\mathcal L\to\mathcal K\) 满足 \(u_{x_1}T_c^{\mathcal L}=T_c^{\mathcal K}u_{x_0}\)，所以换基点同时保持态射以及同构的对应。
\end{proposition}
\begin{proof}
	\(\psi_c\) 是\cref{g2:prop:basepoint-change}中从 \(x_0\) 到 \(x_1\) 的换基点同构的逆。直接展开可见
	\[
		T_{c*h*\bar c}=T_c^{-1}T_hT_c.
	\]
	两边左复合 \(T_c\) 就得到\eqref{g2:eq:basepoint-module-transport}。又因 \(c'\simeq(c'*\bar c)*c\)，运输公式给出 \(T_{c'}=T_cT_g\)。将 \(c'\simeq g*c\) 及其反向道路代入 \(c'*h*\bar c'\)，得到 \(\psi_{c'}(h)=g\psi_c(h)g^{-1}\)。最后的等式就是态射沿 \(c\) 的自然性方形。因此道路选择造成的两项变化彼此匹配，不会改变所描述的局部系统。
\end{proof}

若再接一条 \(d:x_1\to x_2\)，则 \(T_{c*d}=T_dT_c\)，同时 \(\psi_{c*d}=\psi_c\psi_d\)。这说明连续两次换基点与一次沿拼接道路换基点相容。纤维间的比较随连接道路而定；即使基本群交换，非恒等的环路运输仍会改变这个比较。

\begin{proposition}[整体截面与平凡性]\label{g2:prop:parallel-sections}
	设 \(X\) 非空、连通、局部道路连通且半局部单连通，\(\mathcal L\) 为 \(R\) 模局部系统，\(\mathscr F\) 为相应的局部常值层。固定 \(x_0\)，记 \(M=\mathcal L_{x_0}\)、\(\pi=\pi_1(X,x_0)\)。在 \(x_0\) 求值给出自然同构
	\[
		\mathscr F(X)\cong M^\pi,
		\qquad M^\pi=\{m\in M:m\cdot g=m\text{ 对每个 }g\in\pi\}.
	\]
	局部系统 \(\mathcal L\) 平凡，当且仅当每个基点环路作用均为恒等。
\end{proposition}
\begin{proof}
	把整体截面看成纤维覆盖的连续截面。它沿任何道路的限制都是从给定初值出发的提升，由提升唯一性，其终值必为初值的运输；所以基点值确定整个截面，而且被每条环路固定。反过来，给定 \(m\in M^\pi\)，对 \(x\in X\) 选 \(\gamma:x_0\to x\)，令 \(s(x)=T_\gamma m\)。若改用 \(\eta\)，则 \(g=[\gamma*\bar\eta]\) 满足 \([\gamma]=[g*\eta]\)，于是 \(T_\gamma m=T_\eta(m\cdot g)=T_\eta m\)，定义与选择无关。在局部平凡化邻域内，它正是沿小道路运输同一个值所得的截面，故连续。两构造互逆且线性。

	常值系统的环路运输均为恒等，同构保持这一性质。若全部环路作用为恒等，上述方法可以同时用于每个 \(m\in M\)，给出 \(\underline M\to\mathcal L\) 的纤维线性同构，并与道路运输交换，故 \(\mathcal L\) 平凡。
\end{proof}

例如，\(\mathbb Z_-\) 的整体截面只有零，因为 \(m=-m\) 在 \(\mathbb Z\) 中只允许 \(m=0\)。在单连通的局部良好空间上，每个局部系统则平凡；选定一点的纤维后，沿道路运输给出不依赖道路选择的比较。这也是把系数拉到万有覆盖上以后能够使用一个固定模的原因。
% 纲要条目排为编号小节标题，并进入目录与书签；说明段仍保留于纲要注释中。

% !TeX root = ../../Book2.tex
% 纲要标题：### 5.2 扭系数链与上链
% 纲要位置：几何学计划纲要.md，第 1202 行。
% 正文状态：未撰写；本文件目前是框架。

% 纲要内容开始
% 
% * 万有覆盖链复形的群环作用
% * 左右作用及张量、Hom 的约定
% * 局部系数边界和连接映射
% * 同伦不变性与切除
% 
% 纲要内容结束

\section{扭系数链与上链}
\label{g2:sec:0502}

待撰写

\subsection{万有覆盖链复形的群环作用}
\label{g2:subsec:0502:01}
待撰写

\subsection{左右作用及张量、Hom 的约定}
\label{g2:subsec:0502:02}
待撰写

\subsection{局部系数边界和连接映射}
\label{g2:subsec:0502:03}
待撰写

\subsection{同伦不变性与切除}
\label{g2:subsec:0502:04}
待撰写
% 纲要条目排为编号小节标题，并进入目录与书签；说明段仍保留于纲要注释中。

% !TeX root = ../../Book2.tex
% 纲要标题：### 5.3 定向系数和循环模型
% 纲要位置：几何学计划纲要.md，第 1209 行。
% 正文状态：未撰写；本文件目前是框架。

% 纲要内容开始
% 
% * 流形局部同调给出的定向系统
% * 圆周单值下的核、余核计算
% * Möbius 型和不可定向曲面的计算
% * 平凡化系数和遗失的单值信息
% 
% 纲要内容结束

\section{定向系数和循环模型}
\label{g2:sec:0503}

待撰写

\subsection{流形局部同调给出的定向系统}
\label{g2:subsec:0503:01}
待撰写

\subsection{圆周单值下的核、余核计算}
\label{g2:subsec:0503:02}
待撰写

\subsection{Möbius 型和不可定向曲面的计算}
\label{g2:subsec:0503:03}
待撰写

\subsection{平凡化系数和遗失的单值信息}
\label{g2:subsec:0503:04}
待撰写
% 纲要条目排为编号小节标题，并进入目录与书签；说明段仍保留于纲要注释中。

% !TeX root = ../../Book2.tex
% 纲要标题：### 5.4 纤维上同调如何变化
% 纲要位置：几何学计划纲要.md，第 1216 行。
% 正文状态：未撰写；本文件目前是框架。

% 纲要内容开始
% 
% * 纤维运输在同调上的作用
% * 映射环面的扭曲
% * 为 Serre 谱序列和障碍理论准备的系数函子
% 
% ---
% 
% 纲要内容结束

\section{纤维上同调如何变化}
\label{g2:sec:0504}

待撰写

\subsection{纤维运输在同调上的作用}
\label{g2:subsec:0504:01}
待撰写

\subsection{映射环面的扭曲}
\label{g2:subsec:0504:02}
待撰写

\subsection{为 Serre 谱序列和障碍理论准备的系数函子}
\label{g2:subsec:0504:03}
待撰写
% 纲要条目排为编号小节标题，并进入目录与书签；说明段仍保留于纲要注释中。


\begin{chaptersummary}
	局部系统给每点一个模，给每条道路类一个运输同构。在连通、局部道路连通且半局部单连通的空间上，选定基点后，全部数据等价于一个带基本群右作用的模；覆盖纤维空间和局部常值层给出另外两种描述。换基点通过一条连接道路比较纤维，换道路时单值和基本群同构同时改变。

	\ProseParagraphBreak

	万有覆盖的链是左群环自由模。将它与右系数模取张量得到扭系数链；取等变同态时，须先把系数通过逆作用改写为左模。几何上，系数放在单形首顶点：删掉首顶点的面要向前运输系数，上链则要把面上的值逆向送回。单形内的运输相容性保证边界平方为零，棱柱和小链同伦则给出同伦不变性与切除。

	\ProseParagraphBreak

	圆周把全部单值集中在一个自同构 \(T\) 上。一次同调和零次上同调取 \(\ker(T-1)\)，零次同调和一次上同调取 \(\operatorname{coker}(T-1)\)。前者要求一个值绕行后不变，后者把绕行前后的差值置零。无边界流形的局部顶次同调组成定向系统；它在不可定向闭曲面上补偿局部方向的反转，使顶次整数类重新出现。

	\ProseParagraphBreak

	局部平凡纤维丛沿道路给出纤维的同伦运输。同调向前推送，上同调则通过逆运输的拉回向前比较。映射环面的正合列把这些系数系统在圆周上的核、余核送入总空间计算。局部系统保留单值，但不包含丛的全部粘合信息；同调中有不变元也不保证存在几何截面。后续谱序列与障碍理论将在这些明确的运输数据上继续组织更高层信息。
\end{chaptersummary}

\BookExercises
除另行说明外取整数系数，路径拼接与左右作用沿用正文约定。计算扭系数时先写出单值及相应复形；由短正合列得到抽象群的直和分解时，说明分裂的依据。

\begin{exercise}\label{g2:ex:constant-presheaf-gluing}
	设 \(A\) 为非零 \(R\) 模，\(X=U\amalg V\)，其中 \(U,V\) 为非空连通开集。给每个非空开集 \(W\subset X\) 指定 \(P(W)=A\)，并把非空开集之间的限制映射设为恒等，另令 \(P(\varnothing)=0\)。证明这些数据不构成层；计算真正的常值层 \(\underline A(X)\)。再证明 \(P\) 与 \(\underline A\) 在每一点的茎同构，解释为何单独知道全部茎的同构类型不能判定粘合性质。
\end{exercise}
\begin{solution}
	取两个不同的元素 \(a,b\in A\)，分别作为 \(U,V\) 上的截面。它们在空交集上的限制同为零，满足相容条件；若能由 \(c\in P(X)=A\) 粘合，则恒等限制同时要求 \(c=a\) 与 \(c=b\)，矛盾。所以 \(P\) 不满足层的粘合条件。

	一个到离散模 \(A\) 的局部常值函数在连通集上必为常值：任意取值的逆像及其补集均开，若两者都非空就分离了该连通集。因此 \(\underline A(X)=A\times A\)，两坐标为在 \(U,V\) 上的值，任意一对都能唯一粘成局部常值函数。

	在一点 \(x\) 附近，\(P\) 的每个代表只是一个 \(a\in A\)，其进一步限制不改变 \(a\)，故 \(P_x\cong A\)。一个局部常值函数在 \(x\) 的某个邻域上恒等于其值，所以 \(\underline A_x\cong A\)。把 \(a\) 送到恒值函数还给出自然映射 \(P\to\underline A\)，它在每个茎上为同构，在 \(X\) 上却只是对角映射 \(A\to A\times A\)。茎记录已有截面的局部相等关系，不能自行保证不同开集上的相容截面已有一个整体代表；后一项正是这里缺少的粘合公理。
\end{solution}

\begin{exercise}\label{g2:ex:unipotent-circle-system}
	设 \(\mathcal L\) 为圆周上的 \(\mathbb Q\) 模局部系统，基点纤维为列向量空间 \(\mathbb Q^2\)，正向生成环路的运输矩阵为
	\[
		J=\begin{pmatrix}1&1\\0&1\end{pmatrix}.
	\]
	对整数 \(k\) 定义 \(f_k:S^1\to S^1\)、\(f_k(z)=z^k\)，其中 \(f_0\) 为常值映射。求 \(f_k^*\mathcal L\) 的单值及整体截面，判断哪些 \(k\) 给出平凡系统，并证明当 \(k\ne0\) 时 \(f_k^*\mathcal L\cong\mathcal L\)。
\end{exercise}
\begin{solution}
	正向环路经 \(f_k\) 绕行 \(k\) 周，故拉回系统的生成运输为 \(J^k\)。写 \(J=I+N\)，有 \(N^2=0\) 及 \((I+N)^{-1}=I-N\)，所以对所有整数 \(k\) 都有
	\[
		J^k=I+kN=\begin{pmatrix}1&k\\0&1\end{pmatrix}.
	\]
	由\cref{g2:prop:parallel-sections}，整体截面对应于 \(\ker(J^k-I)\)。当 \(k=0\) 时这是全部 \(\mathbb Q^2\)，系统平凡；当 \(k\ne0\) 时这是 \(\mathbb Q(1,0)^{\mathsf T}\)，且 \(J^k\ne I\)，故系统不平凡。

	对 \(k\ne0\)，令 \(S=\operatorname{diag}(k,1)\)，则 \(S\) 可逆，且 \(SJ=J^kS\)。这说明 \(S\) 保持原系统与拉回系统的右生成元作用，因此由\cref{g2:thm:local-system-classification}延拓为局部系统同构。拉回改变了指定基下的单值矩阵，但并不一定改变局部系统的同构类。
\end{solution}

\begin{exercise}\label{g2:ex:projective-rank-one-system}
	分类 \(\mathbb{RP}^2\) 上纤维同构于 \(\mathbb Z\) 的整数局部系统。对反极点覆盖 \(p:S^2\to\mathbb{RP}^2\)，计算这些系统的拉回，并比较它们在 \(\mathbb{RP}^2\) 上的整体截面。
\end{exercise}
\begin{solution}
	\(S^2\) 单连通，而 \(p\) 是二叶覆盖，所以它是万有覆盖，覆盖变换为恒等和反极点映射。由\cref{g2:prop:monodromy-deck}，\(\pi_1(\mathbb{RP}^2)\cong\mathbb Z/2\mathbb Z\)；同一计算也见\cref{g2:ex:projective-pi1}。两个空间均是连通流形，具备分类定理要求的局部条件。整数模 \(\mathbb Z\) 的自同构只有乘以 \(1\) 和乘以 \(-1\)，两者的平方均为恒等，所以恰有两种右作用。它们不共轭，因为 \(\operatorname{Aut}_{\mathbb Z}(\mathbb Z)=\{1,-1\}\) 是交换群。因此有且仅有两个同构类，分别由平凡单值与符号单值给出。

	两个系统拉回到 \(S^2\) 后均平凡，因为 \(\pi_1(S^2)=0\)，没有非平凡环路作用。原空间上，平凡系统的整体截面为 \(\mathbb Z\)，符号系统的整体截面为 \(\{m\in\mathbb Z:m=-m\}=0\)。所以在万有覆盖上平凡，并不意味着原局部系统平凡；下降回底空间时仍需记住覆盖变换怎样作用于系数。
\end{solution}

\begin{exercise}\label{g2:ex:noncommuting-transport}
	令 \(X=S^1\vee S^1\)，取两条正向环路 \(a,b\) 为自由生成元。设 \(\mathbb Q^2\) 上的生成运输为
	\[
		T_a=A=\begin{pmatrix}1&1\\0&1\end{pmatrix},\qquad
		T_b=B=\begin{pmatrix}1&0\\1&1\end{pmatrix}.
	\]
	说明这些数据确定一个局部系统，计算 \(T_{a*b}\) 和 \(T_{a*b*\bar a*\bar b}\)。若换基点的两条道路 \(c,c'\) 满足 \([c'*\bar c]=[a*b]\)，写出 \(T_{c'}\) 与 \(T_c\) 的关系，并说明为何不能把 \(T_aT_b\) 直接当成 \(T_{a*b}\)。
\end{exercise}
\begin{solution}
	\cref{g2:thm:graph-free-group}给出 \(\pi_1(X)=\langle a,b\rangle\)，没有关系需要检验。可先定义左表示 \(\rho(a)=A^{-1}\)、\(\rho(b)=B^{-1}\)，由自由群性质延拓到整个基本群，再令 \(T_g=\rho(g)^{-1}\)。它满足 \(T_{gh}=T_hT_g\)，所以由\cref{g2:thm:local-system-classification}确定所求局部系统。

	按行进顺序复合运输，得到
	\[
		T_{a*b}=BA=\begin{pmatrix}1&1\\1&2\end{pmatrix},\qquad
		T_{a*b*\bar a*\bar b}=B^{-1}A^{-1}BA
		=\begin{pmatrix}0&-1\\1&3\end{pmatrix}.
	\]
	这里可先计算 \(BA\)，左乘 \(A^{-1}\) 得 \(\left(\begin{smallmatrix}0&-1\\1&2\end{smallmatrix}\right)\)，再左乘 \(B^{-1}\) 得所示矩阵。换基点公式给出 \(T_{c'}=T_cT_{a*b}=T_cBA\)。函数复合先应用右边矩阵，故先沿 \(a\) 运输再沿 \(b\) 运输必须是 \(BA\)。由于 \(AB=\left(\begin{smallmatrix}2&1\\1&1\end{smallmatrix}\right)\ne BA\)，在这个例子中倒转顺序会实质改变答案，而不只是改变记号。
\end{solution}

\begin{exercise}\label{g2:ex:twisted-interval-pair}
	设 \(I=[0,1]\)，\(A=\{0,1\}\)，\(\mathcal L\) 为 \(I\) 上的 \(R\) 模局部系统，\(M=\mathcal L_0\)，\(T:\mathcal L_0\to\mathcal L_1\) 为沿正向区间的运输。计算 \(H_n(I,A;\mathcal L)\) 和 \(H^n(I,A;\mathcal L)\)，并在不把两个端点纤维预先视为同一个模的情况下写出同调连接映射。
\end{exercise}
\begin{solution}
	沿区间内唯一的同端点道路类运输，可以把整个系统平凡化为 \(M\)。区间和端点都是 CW 空间，所以\cref{g2:thm:twisted-cellular-comparison}适用。相对胞腔复形只保留一条一胞腔，系数为 \(M\)，边界为零，故
	\[
		H_n(I,A;\mathcal L)\cong H^n(I,A;\mathcal L)\cong
		\begin{cases}M,&n=1,\\0,&n\ne1.\end{cases}
	\]
	这里上同调的一次识别也使用从首顶点出发的平凡化；相对上链在端点上为零，故只有一维项。若用正向一胞腔把一次相对同调识别为 \(M\)，\eqref{g2:eq:twisted-boundary}给出
	\[
		\partial_*:M\longrightarrow H_0(A;\mathcal L)=\mathcal L_0\oplus\mathcal L_1,
		\qquad m\longmapsto(-m,Tm).
	\]
	两个端点的系数不能直接相减；它们是落在两个直和项中的元素。映入 \(H_0(I;\mathcal L)\) 后，正向路径的边界恰把 \(m\) 与 \(Tm\) 识别起来，所以这对元素的像为零，符合长正合列。
\end{solution}

\begin{exercise}\label{g2:ex:two-loop-local-matrices}
	设 \(X=S^1\vee S^1\)，两条正向环路记为 \(a,b\)。在基点纤维 \(M=\mathbb Z^2\) 上规定运输
	\[
		T_a(x,y)=(-x,y),\qquad T_b(x,y)=(x,-y).
	\]
	计算所得局部系统 \(\mathcal L\) 的全部同调与上同调群，并分别给出 \(H_1\) 的生成元和 \(H^1\) 的自由部分、挠部分的代表。
\end{exercise}
\begin{solution}
	楔和有一个零胞腔、两条一胞腔，没有更高胞腔。两条提升边的边界分别为 \((a-1)\widetilde v\)、\((b-1)\widetilde v\)。同调复形为
	\[
		0\longrightarrow M\oplus M\xrightarrow{\mathrm d}M\longrightarrow0,
		\qquad \mathrm d((x_1,x_2),(y_1,y_2))=(-2x_1,-2y_2).
	\]
	因此核由“第一条边的系数 \((0,1)\)”和“第二条边的系数 \((1,0)\)”生成，像为 \(2\mathbb Z\oplus2\mathbb Z\)，得到
	\[
		H_1(X;\mathcal L)=\mathbb Z^2,\qquad
		H_0(X;\mathcal L)=(\mathbb Z/2)^2,
		\qquad H_n(X;\mathcal L)=0\quad(n\ge2).
	\]
	由于两个运输都是自身的逆，上链复形为
	\[
		0\longrightarrow M\xrightarrow{\delta}M\oplus M\longrightarrow0,
		\qquad\delta(u,v)=((-2u,0),(0,-2v)).
	\]
	其核为零；商群中第一条边的第二坐标、第二条边的第一坐标给出两个自由生成元，而第一条边的第一坐标、第二条边的第二坐标分别只按偶数关系取商，给出两个二阶生成元。因此
	\[
		H^0(X;\mathcal L)=0,\qquad
		H^1(X;\mathcal L)=\mathbb Z^2\oplus(\mathbb Z/2)^2,
		\qquad H^n(X;\mathcal L)=0\quad(n\ge2).
	\]
	这里零次同调与零次上同调已不同：前者允许把运输差值置零，后者要求一个值同时被两个运输固定。
\end{solution}

\begin{exercise}\label{g2:ex:regular-local-coefficients-circle}
	令 \(\Gamma=\mathbb Z[t,t^{-1}]\)，在圆周上取基点纤维为右正则模 \(M=\Gamma\)、正向单值为 \(m\mapsto mt\) 的局部系统 \(\mathcal R\)。计算 \(H_*(S^1;\mathcal R)\) 与 \(H^*(S^1;\mathcal R)\)，并说明为什么后一组群不能直接解释为万有覆盖 \(\mathbb R\) 的普通整数上同调。
\end{exercise}
\begin{solution}
	胞腔链只有一次和零次，边界为乘 \(t-1\)。Laurent 多项式环是整环，故这个乘法单射；取值 \(t=1\) 的同态满射到 \(\mathbb Z\)，核为 \((t-1)\Gamma\)。核的断言可将 Laurent 多项式乘以足够大的 \(t\) 次幂，化为普通多项式的因式分解；\(t\) 在 \(\Gamma\) 中可逆，故不会改变该理想。因此
	\[
		H_0(S^1;\mathcal R)=\mathbb Z,\qquad H_n(S^1;\mathcal R)=0\quad(n>0).
	\]
	上链余边界为乘 \(t^{-1}-1=-t^{-1}(t-1)\)。因 \(-t^{-1}\) 是单位，这个映射仍单射且像仍为 \((t-1)\Gamma\)，所以
	\[
		H^1(S^1;\mathcal R)=\mathbb Z,\qquad H^n(S^1;\mathcal R)=0\quad(n\ne1).
	\]
	同调模型中 \(\Gamma\otimes_\Gamma C_*(\mathbb R;\mathbb Z)\cong C_*(\mathbb R;\mathbb Z)\)，所以同调确实恢复万有覆盖的普通整数同调。上链模型却是 \(\operatorname{Hom}_\Gamma(C_*(\mathbb R;\mathbb Z),\Gamma^-)\)，要求等变性且取值于群环；普通上链为 \(\operatorname{Hom}_{\mathbb Z}(C_*(\mathbb R;\mathbb Z),\mathbb Z)\)，两者的对象、取值和等变条件都不同。\(\mathbb R\) 可缩，其普通整数上同调仅在零次为 \(\mathbb Z\)，与上述结果并不矛盾。
\end{solution}

\begin{exercise}\label{g2:ex:sign-coefficient-connecting}
	在圆周上令 \(\mathcal Z_-\) 为单值 \(-1\) 的整数局部系统，\(\mathbb F_2\) 为常值系统。模 \(2\) 把 \(-1\) 化为 \(1\)，故有局部系统短正合列
	\[
		0\longrightarrow\mathcal Z_-\xrightarrow{\times2}\mathcal Z_-\longrightarrow\mathbb F_2\longrightarrow0.
	\]
	证明它给出同调及上同调的系数长正合列，并计算连接同态
	\[
		H_1(S^1;\mathbb F_2)\longrightarrow H_0(S^1;\mathcal Z_-),\qquad
		H^0(S^1;\mathbb F_2)\longrightarrow H^1(S^1;\mathcal Z_-).
	\]
\end{exercise}
\begin{solution}
	局部系数链逐单形取纤维的直和，所以上述系数序列给出短正合链复形。上链逐单形取纤维的直积；每个模 \(2\) 值可选整数代表，因而末项仍满射，也给出短正合上链复形。运输相容性使系数映射与边界、余边界交换。于是\cref{g2:thm:homology-LES}的提升构造及其上链形式适用：先把模 \(2\) 的圈或上圈提升为整数链或上链，再取边界，其值必为偶数，最后除以 \(2\) 得到连接类。

	用一顶点一边的胞腔模型，符号系数的链边界和上链余边界均为乘 \(-2\)，模 \(2\) 系数的两种微分均为零。因此
	\[
		H_1(S^1;\mathbb F_2)=H^0(S^1;\mathbb F_2)=\mathbb F_2,
		\qquad H_0(S^1;\mathcal Z_-)=H^1(S^1;\mathcal Z_-)=\mathbb Z/2.
	\]
	一次模 \(2\) 链生成元可提升为系数 \(1\) 的整数边，其边界为 \(-2\)，除以 \(2\) 得到 \(-1\)，代表 \(\mathbb Z/2\) 的非零元。零次模 \(2\) 上链生成元提升为顶点值 \(1\)，余边界同样为 \(-2\)，所以第二个连接映射也把 \(1\) 送到 \([-1]\)。两映射都是同构。胞腔比较可在选取 \(J\) 后与所有系数映射同时使用，故胞腔中算出的连接类就是奇异系数正合列中的连接类。
\end{solution}

\begin{exercise}\label{g2:ex:integer-unipotent-circle}
	取整数 \(q\geq2\)，令 \(A=\mathbb Z^2\)，单值自同构为
	\[
		T(x,y)=(x+qy,y).
	\]
	计算圆周上的 \(H_k(S^1;\mathcal L_T)\) 和 \(H^k(S^1;\mathcal L_T)\)，并给出核、余核中所用的生成元。再将系数逐纤维改成 \(\mathbb Q\)，比较结果。说明这次出现的挠元怎样来自单值。
\end{exercise}
\begin{solution}
	记标准基为 \(e_1,e_2\)，则 \((T-1)(x,y)=(qy,0)\)，所以
	\[
		\ker(T-1)=\mathbb Ze_1,\qquad
		(T-1)A=q\mathbb Ze_1,\qquad
		A/(T-1)A=(\mathbb Z/q)[e_1]\oplus\mathbb Z[e_2].
	\]
	由\cref{g2:prop:circle-local-coefficients}，\(H_1=H^0=\mathbb Ze_1\)，\(H_0=H^1=(\mathbb Z/q)[e_1]\oplus\mathbb Z[e_2]\)，更高次数为零。上链微分用 \(T^{-1}-1\)，其公式为 \((x,y)\mapsto(-qy,0)\)，确实具有相同的核与像。

	改为有理数系数后，像为整个 \(\mathbb Qe_1\)，核为 \(\mathbb Qe_1\)，余核为 \(\mathbb Q[e_2]\)。因此零次、一次同调与上同调都同构于 \(\mathbb Q\)，更高次数为零。整数挠元 \([e_1]\) 的阶恰为 \(q\)：绕行把 \(e_2\) 变成 \(e_2+qe_1\)，零次同调要求两者相等，留下关系 \(qe_1=0\)，却没有关系 \(e_1=0\)。这说明扭系数计算中的挠元不仅能来自空间本身的整数同调，也能来自单值矩阵的整数像。
\end{solution}

\begin{exercise}\label{g2:ex:surface-other-sign-systems}
	设 \(h\geq2\)，在 \(N_h\) 的标准边 \(a_1,\ldots,a_h\) 上指定符号 \(\varepsilon_i\in\{1,-1\}\)，且至少一个符号为 \(1\)，至少一个为 \( -1\)。证明这些符号定义一个秩一整数局部系统 \(\mathcal L_\varepsilon\)，并计算所有 \(H_k(N_h;\mathcal L_\varepsilon)\)。将所得结果与常值系统及定向系统比较。
\end{exercise}
\begin{solution}
	关系 \(a_1^2\cdots a_h^2=1\) 在符号作用下变成 \(\varepsilon_1^2\cdots\varepsilon_h^2=1\)，所以各生成元的指定作用满足群关系，确实给出一个右 \(\mathbb Z[\pi_1(N_h)]\) 模，进而由\cref{g2:thm:local-system-classification}得到局部系统。

	令 \(I=\{i:\varepsilon_i=-1\}\)、\(J=\{i:\varepsilon_i=1\}\)，两集合均非空。提升多边形公式\eqref{g2:eq:lifted-word-boundary}中，各前缀 \(Q_i=a_1^2\cdots a_{i-1}^2\) 的作用均为 \(1\)，故胞腔复形为
	\[
		0\longrightarrow\mathbb Z
		\xrightarrow{\ 1\mapsto2\sum_{j\in J}e_j\ }\mathbb Z^h
		\xrightarrow{\ (x_i)\mapsto-2\sum_{i\in I}x_i\ }\mathbb Z
		\longrightarrow0.
	\]
	第一映射因 \(J\ne\varnothing\) 单射，所以 \(H_2=0\)；末端映射的像因 \(I\ne\varnothing\) 恰为 \(2\mathbb Z\)，所以 \(H_0=\mathbb Z/2\)。

	固定 \(i_0\in I\)，中间核有一组整数基
	\[
		e_i-e_{i_0}\quad(i\in I\setminus\{i_0\}),\qquad e_j\quad(j\in J).
	\]
	令 \(u_J=\sum_{j\in J}e_j\)。在后一组基中，选定一个 \(j_0\in J\)，以 \(u_J\) 代替 \(e_{j_0}\) 仍为整数基，因为 \(e_{j_0}=u_J-\sum_{j\ne j_0}e_j\)。二维边界仅给出关系 \(2u_J=0\)，于是
	\[
		H_1(N_h;\mathcal L_\varepsilon)\cong\mathbb Z^{h-2}\oplus\mathbb Z/2.
	\]
	高于二次没有胞腔，故其群为零。所有符号为 \(1\) 时是常值系统，结果为 \(H_0=\mathbb Z\)、\(H_1=\mathbb Z^{h-1}\oplus\mathbb Z/2\)、\(H_2=0\)；所有符号为 \( -1\) 时是定向系统，结果为 \(H_0=\mathbb Z/2\)、\(H_1=\mathbb Z^{h-1}\)、\(H_2=\mathbb Z\)。因此不是任意非平凡秩一系统都能恢复顶次整数类，定向系统的符号与几何附着相匹配才产生这一效果。
\end{solution}

\begin{exercise}\label{g2:ex:moebius-relative-orientation}
	令 \(B=(\mathbb R\times[-1,1])/((s,u)\sim(s+1,-u))\) 为闭 Möbius 带，\(r:B\to S^1\) 为核心收缩，取 \(\mathcal L=r^*\mathbb Z_{-}\)。证明 \(\mathcal L|_{\partial B}\) 为常值系统，并计算 \(H_k(B,\partial B;\mathcal L)\)。在长正合列中明确写出 \(H_0(\partial B;\mathcal L)\to H_0(B;\mathcal L)\)。
\end{exercise}
\begin{solution}
	边界是一个圆周：从 \([0,1]\) 出发沿上边走到 \([1,1]=[0,-1]\)，再沿下边走到 \([1,-1]=[0,1]\)，恰绕完边界一次。投到核心时，这条路径绕行两周。因此边界生成元的单值为 \((-1)^2=1\)，边界上的系统为常值整数系统；选一个边界基点处的生成元，就得到其平凡化。

	由核心收缩和\cref{g2:prop:circle-local-coefficients}，
	\[
		H_1(B;\mathcal L)=0,\qquad H_0(B;\mathcal L)=\mathbb Z/2,\qquad
		H_1(\partial B;\mathcal L)=H_0(\partial B;\mathcal L)=\mathbb Z.
	\]
	沿一条从边界基点到核心基点的路径识别纤维，包含诱导的零次映射将整数 \(m\) 送到其在 \(\mathbb Z/2\) 中的类。换一条识别路径至多反号，在模二余核中不改变该映射。故空间对的长正合列在低次成为
	\[
		\begin{aligned}
			0\longrightarrow H_2(B,\partial B;\mathcal L)\longrightarrow\mathbb Z\longrightarrow0,\\
			0\longrightarrow H_1(B,\partial B;\mathcal L)\longrightarrow\mathbb Z\xrightarrow{\bmod\,2}\mathbb Z/2
			\longrightarrow H_0(B,\partial B;\mathcal L)\longrightarrow0.
		\end{aligned}
	\]
	正合性给出 \(H_2(B,\partial B;\mathcal L)\cong\mathbb Z\)、\(H_1(B,\partial B;\mathcal L)\cong2\mathbb Z\cong\mathbb Z\)、\(H_0(B,\partial B;\mathcal L)=0\)。高次由两端的消失同样为零。这里相对二维类的边界生成 \(H_1(\partial B;\mathcal L)\)，而相对一次类到边界零次群的像是 \(2\mathbb Z\)；两条连接映射的整数倍数不同。
\end{solution}

\begin{exercise}\label{g2:ex:product-orientation-system}
	设 \(M,N\) 分别为连通、Hausdorff、第二可数、无边界的 \(m,n\) 维拓扑流形，\(m,n\geq1\)。固定乘积中先列 \(M\) 坐标、后列 \(N\) 坐标的次序。证明相对外积给出局部系统同构
	\[
		\mathcal O_{M\times N}\cong
		\operatorname{pr}_M^*\mathcal O_M\otimes_{\mathbb Z}\operatorname{pr}_N^*\mathcal O_N.
	\]
	推出 \(M\times N\) 可定向当且仅当 \(M,N\) 都可定向，并解释为何两个不可定向流形相乘不会自动使定向反转抵消。
\end{exercise}
\begin{solution}
	在点 \((x,y)\) 处，两个被删去的单点在 Hausdorff 流形中闭，故其补空间开，且
	\[
		(M\setminus\{x\})\times N\ \cup\ M\times(N\setminus\{y\})
		=(M\times N)\setminus\{(x,y)\}.
	\]
	因此\cref{g2:prop:cross-cup-evaluation}证明中构造的相对同调外积适用，得到
	\[
		\mathcal O_{M,x}\otimes_{\mathbb Z}\mathcal O_{N,y}
		\longrightarrow\mathcal O_{M\times N,(x,y)}.
	\]
	切除到两个坐标球后，\cref{g2:prop:local-cross-normalization}的乘积基本链证明说明：两个局部生成元的外积仍为乘积邻域的生成元，且坐标顺序固定其正号。两端均为无限循环群，因此此映射是同构。

	这些同构还必须保持运输。分别在两个坐标球内用共同相对群延拓生成元，其外积在乘积邻域内仍为同一个相对类在各点的像；相对外积的自然性保证这些延拓与各点同构交换。分段任意乘积路径并在每段使用这种乘积邻域，就得运输相容性。于是逐点同构组成局部系统同构。

	以 \((x_0,y_0)\) 为基点，乘积闭路 \((\alpha,\beta)\) 的作用为
	\[
		w_{M\times N}([\alpha],[\beta])=w_M([\alpha])\,w_N([\beta]).
	\]
	若两个因子都可定向，右侧恒为 \(1\)，乘积可定向。反之，只在第一因子走闭路、第二因子保持不动，可知乘积可定向迫使 \(w_M\equiv1\)；交换两因子的角色得到 \(w_N\equiv1\)。两个反定向闭路同时走时，符号确实相乘为 \(1\)，但乘积空间还包含只在一个因子中绕行的闭路，后者的负号没有抵消。
\end{solution}

\begin{exercise}\label{g2:ex:torus-unipotent-mapping-torus}
	令 \(F=\mathbb R^2/\mathbb Z^2\)，矩阵
	\[
		A=\begin{pmatrix}1&2\\0&1\end{pmatrix}
	\]
	诱导同胚 \(f:F\to F\)。计算映射环面 \(E_f\) 的全部整数同调与上同调群。说明为何纤维和基底的整数同调都无挠，而总空间仍能出现挠元，并判断此丛是否具有截面。
\end{exercise}
\begin{solution}
	矩阵 \(A\) 及其逆都是整数矩阵，故确实在商环面上诱导互逆的连续映射。环面同调为 \(H_0=H_2=\mathbb Z\)、\(H_1=\mathbb Z^2\)。\(f_*\) 在一次同调上的矩阵是 \(A\)，在零次为恒等；在二次为乘 \(\det A=1\)，可由\cref{g2:prop:torus-cohomology-ring}的两个一次类相乘计算得到。于是
	\[
		(A-I)(x,y)=(2y,0),\qquad
		\ker(A-I)=\mathbb Ze_1,\qquad
		\operatorname{coker}(A-I)=(\mathbb Z/2)[e_1]\oplus\mathbb Z[e_2].
	\]
	将这些群放入\cref{g2:thm:mapping-torus-sequences}，各短正合列的商群均为自由阿贝尔群，故可选择基的提升使之分裂。得到
	\[
		H_k(E_f;\mathbb Z)\cong
		\begin{cases}
			\mathbb Z,&k=0,3,\\
			\mathbb Z^2\oplus\mathbb Z/2,&k=1,\\
			\mathbb Z^2,&k=2,\\
			0,&k>3.
		\end{cases}
	\]
	上同调中，\(f^*\) 在对偶一次基上的矩阵为 \(A^{\mathsf T}\)，所以 \((A^{\mathsf T}-I)(u,v)=(0,2u)\)。其核为 \(\mathbb Ze_2^*\)，余核为 \(\mathbb Z[e_1^*]\oplus(\mathbb Z/2)[e_2^*]\)；在零次和二次仍为恒等。上同调短正合列的商群同样自由，故
	\[
		H^k(E_f;\mathbb Z)\cong
		\begin{cases}
			\mathbb Z,&k=0,3,\\
			\mathbb Z^2,&k=1,\\
			\mathbb Z^2\oplus\mathbb Z/2,&k=2,\\
			0,&k>3.
		\end{cases}
	\]
	这里同调挠元来自单值关系 \(f_*e_2-e_2=2e_1\)，上同调挠元来自对偶关系；纤维群本身自由，并不保证运输差值的像是直和项。原点是 \(f\) 的固定点，故 \(e^{2\pi\mathrm{i}t}\mapsto[0,t]\) 给出截面。截面只提供一个基底方向，不会消除纤维单值造成的挠元。
\end{solution}

\begin{exercise}\label{g2:ex:rotation-mapping-torus}
	取 \(\theta\in\mathbb R\setminus\mathbb Z\)，令 \(f:S^1\to S^1\)、\(f(z)=e^{2\pi\mathrm{i}\theta}z\)。证明 \(f\) 没有固定点，但 \(E_f\to S^1\) 有截面。进一步构造一个覆盖基底恒等映射的同胚 \(E_f\cong S^1\times S^1\)，并写出它给出的截面在柱坐标中的道路。
\end{exercise}
\begin{solution}
	若 \(f(z)=z\)，因 \(z\ne0\) 得 \(e^{2\pi\mathrm{i}\theta}=1\)，与 \(\theta\notin\mathbb Z\) 矛盾。定义
	\[
		\Phi:E_f\longrightarrow S^1\times S^1,
		\qquad \Phi([z,t])=(e^{2\pi\mathrm{i}\theta t}z,e^{2\pi\mathrm{i}t}).
	\]
	在 \(t=1\) 处的值为 \((e^{2\pi\mathrm{i}\theta}z,1)\)，与 \([f(z),0]\) 的值一致，故此连续公式下降到商空间。给定 \((w,e^{2\pi\mathrm{i}t})\)，取 \(0\leq t<1\)，其原像为 \([e^{-2\pi\mathrm{i}\theta t}w,t]\)，因此 \(\Phi\) 双射。定义域是紧空间的商，靶为 Hausdorff 空间，所以连续双射为同胚；第二坐标正是原投影。

	乘积中的截面 \(b\mapsto(1,b)\) 对应于道路 \(c(t)=e^{-2\pi\mathrm{i}\theta t}\)。它满足 \(c(0)=1=f(c(1))\)，因此 \(t\mapsto[c(t),t]\) 下降为截面。这个例子把“单值同胚有固定点”和“丛有截面”区分开：截面可以沿纤维移动，不必由某个固定纤维点给出。
\end{solution}

\begin{exercise}\label{g2:ex:permutation-mapping-torus}
	设 \(F\) 是含 \(d\geq1\) 个点的离散空间，\(f\) 是其一个置换，循环分解含 \(r\) 个循环，长度为 \(\ell_1,\ldots,\ell_r\)。计算 \(E_f\) 的同调与上同调，给出其道路分支，并判断投影 \(E_f\to S^1\) 何时有截面。解释为何 \(H_0(F;\mathbb Z)\) 总有非零单值不变元，却不总有截面。
\end{exercise}
\begin{solution}
	每个点贡献一条区间，末端按置换接到下一点的区间首端。同一循环的 \(\ell_j\) 条区间依次粘成一个圆周，不同循环之间不相交。因此 \(E_f\) 是 \(r\) 个圆周的不交并，第 \(j\) 个分支到基底的投影绕行 \(\ell_j\) 周。直接得到
	\[
		H_0(E_f;\mathbb Z)=H_1(E_f;\mathbb Z)=\mathbb Z^r,
		\qquad H^0(E_f;\mathbb Z)=H^1(E_f;\mathbb Z)=\mathbb Z^r,
	\]
	其余非负次数的群为零。也可从正合列计算：\(H_0(F)=\mathbb Z^d\)，\(f_*\) 置换点基；\(\ker(f_*-1)\) 由各循环内全部点的和生成，\(\operatorname{coker}(f_*-1)\) 则将同一循环中的点基识别，二者均为 \(\mathbb Z^r\)。高次纤维同调为零，故没有其他贡献；上同调用对偶置换得到相同结果。

	由\cref{g2:prop:mapping-torus-section}，截面存在恰当某个纤维道路分支被 \(f\) 固定。离散空间的道路分支是单点，所以条件正是某个 \(\ell_j=1\)。循环中的点之和总是非零不变元，但当循环长度大于一时，它是多个点的形式和，不是一个能单独返回自身的纤维点。这正是同调不变元与几何截面的差别。
\end{solution}
